\section{Hitung Tangan Paket Cirebon-1}
\label{app:hitung}

Hitung tangan disusun independen dari kode aplikasi untuk preset Paket Cirebon-1: nilai ekspor 65\%, kuota 1.500 MW per tahun, gedung pemerintah 100 MW per tahun, 5\% subsidi dialihkan, Cirebon-1 pensiun 2035, dan 25\% kompensasi ditanggung PLN. Selisih dengan hasil aplikasi di bawah 0,01\% untuk seluruh komponen.

\begin{table}[H]
\centering
\caption{Hitung tangan neraca PLN Paket Cirebon-1 terhadap baseline}
\small
\begin{tabular}{lrrr}
\toprule
Komponen (Rp T, kecuali disebut) & 2030 & 2035 & 2040 \\
\midrule
Kapasitas PLTS atap (MWp) & 5.937 & 12.525 & 19.112 \\
Pendapatan PLN hilang & -4,80 & -10,54 & -16,29 \\
Kompensasi ekspor PLTS & -2,87 & -6,44 & -10,01 \\
Biaya energi dihemat & 3,03 & 6,75 & 10,55 \\
Pembayaran kapasitas dihemat & 0,00 & 1,91 & 1,91 \\
Porsi PLN kompensasi pensiun & 0,00 & -0,17 & -0,17 \\
Pembangkit pengganti & 0,00 & -1,31 & -1,31 \\
Neraca PLN & -4,64 & -9,80 & -15,32 \\
\bottomrule
\end{tabular}

\end{table}

Langkah perhitungan:
\begin{enumerate}
  \item PVOUT rata-rata tertimbang penjualan listrik PLN per provinsi: 1.423,5 kWh/kWp per tahun.
  \item Balik modal rumah tangga non subsidi tanpa ekspor: Rp15 juta / (1.423,5 $\times$ 1.444,70 $\times$ 0,4) = 18,2 tahun. Dengan ekspor 65\%: 9,2 tahun.
  \item Tambahan pasar = tren $\times$ porsi segmen $\times$ exp($-$0,3 $\times$ 9,2) / exp($-$0,3 $\times$ 18,2), naik bertahap selama 4 tahun dan dibatasi kuota.
  \item Biaya marjinal armada PLTU: Rp494 per kWh dari harga DMO US\$70 per ton dan heat rate GEM.
  \item Cirebon-1 pada 2035: pembayaran kapasitas 660 MW $\times$ US\$175/kW $\times$ Rp16.500 = Rp1,91 triliun. Porsi PLN kompensasi 25\% $\times$ Rp1,045 miliar $\times$ 660 MW = Rp0,17 triliun. Pembangkit pengganti 660 MW $\times$ US\$120/kW $\times$ Rp16.500 = Rp1,31 triliun.
\end{enumerate}

\section{Tabel Asumsi}
\label{app:asumsi}

Status ``terverifikasi'' berarti angka dicocokkan dengan dokumen sumber. Kode sumber mengacu pada daftar pustaka (lihat juga halaman Sumber di aplikasi).

{\small
\begin{longtable}{p{0.33\textwidth} p{0.2\textwidth} p{0.26\textwidth} p{0.11\textwidth}}
\toprule
Parameter & Nilai & Sumber & Status \\
\midrule
\endhead
Tahun dasar & 2025 tahun & \texttt{-\allowbreak{}} & asumsi \\
Beban puncak tertinggi Jawa-Madura-Bali & 28.077 MW & \texttt{B-\allowbreak{}MATAHARI-\allowbreak{}2024} & terverifikasi \\
Reserve margin Jawa-Madura-Bali 2025 & 0,375 fraksi & \texttt{B-\allowbreak{}PETROMINDO-\allowbreak{}RESERVE} & terverifikasi \\
Pertumbuhan beban dan penjualan & 0,0445 per tahun & \texttt{B-\allowbreak{}PETROMINDO-\allowbreak{}RESERVE} & terverifikasi \\
Tambahan daya mampu 2026-2028 & 1.000 MW/tahun & \texttt{B-\allowbreak{}PETROMINDO-\allowbreak{}RESERVE} & asumsi \\
Batas bawah reserve margin yang direkomendasikan & 0,33 fraksi & \texttt{B-\allowbreak{}PETROMINDO-\allowbreak{}RESERVE} & terverifikasi \\
Penjualan listrik PLN 2025 & 317,69 TWh & \texttt{B-\allowbreak{}CNBC-\allowbreak{}SALES-\allowbreak{}2025} & terverifikasi \\
Porsi penjualan Jawa-Bali & 0,695 fraksi & \texttt{D-\allowbreak{}PLN-\allowbreak{}STAT-\allowbreak{}2025} & terverifikasi \\
Biaya pokok penyediaan 2026 & 1.874 Rp/kWh & \texttt{B-\allowbreak{}KONTAN-\allowbreak{}SUBSIDI-\allowbreak{}2026} & terverifikasi \\
Kurs & 16.500 Rp/US\$ & \texttt{B-\allowbreak{}KONTAN-\allowbreak{}SUBSIDI-\allowbreak{}2026} & terverifikasi \\
Tingkat diskonto NPV & 0,08 per tahun & \texttt{assumption} & asumsi \\
Porsi kerugian PLN yang ditutup APBN & 1 fraksi & \texttt{assumption} & asumsi \\
Mulai tahun PLN menambah pembangkit untuk menjaga target cadangan & 2029 tahun & \texttt{B-\allowbreak{}LISTRIKID-\allowbreak{}RUPTL-\allowbreak{}JAMALI} & asumsi \\
Harga batu bara DMO untuk PLN & 70 US\$/ton & \texttt{B-\allowbreak{}KONTAN-\allowbreak{}DMO-\allowbreak{}2026} & terverifikasi \\
Kalori acuan harga DMO & 6.322 kcal/kg & \texttt{B-\allowbreak{}KONTAN-\allowbreak{}DMO-\allowbreak{}2026} & perlu cek \\
Biaya variabel non bahan bakar PLTU & 50 Rp/kWh & \texttt{assumption} & asumsi \\
Pembayaran kapasitas IPP (komponen A) & 175 US\$/kW/tahun & \texttt{P-\allowbreak{}IEEFA-\allowbreak{}2017} & asumsi \\
Biaya tetap O\&M PLTU milik PLN & 40 US\$/kW/tahun & \texttt{assumption} & asumsi \\
Nilai buku PLTU milik PLN tahun 2022 & 762.000 US\$/MW & \texttt{B-\allowbreak{}PELABUHANRATU-\allowbreak{}VALUASI} & perlu cek \\
Kompensasi pensiun dini, estimasi skema ETM & 1,04 miliar Rp/MW/tahun dipercepat & \texttt{B-\allowbreak{}REUTERS-\allowbreak{}CIREBON-\allowbreak{}2022} & terverifikasi \\
Kompensasi pensiun dini, klaim PLN & 18,18 miliar Rp/MW/tahun dipercepat & \texttt{B-\allowbreak{}PETROMINDO-\allowbreak{}CIREBON-\allowbreak{}BATAL} & terverifikasi \\
Pekerja langsung per MW & 0,4545 orang/MW & \texttt{B-\allowbreak{}ECOBUSINESS-\allowbreak{}CIREBON} & terverifikasi \\
Biaya kapasitas pembangkit pengganti & 120 US\$/kW/tahun & \texttt{assumption} & asumsi \\
Kematian dini rata-rata armada PLTU & 258,6 kematian/GW/tahun & \texttt{P-\allowbreak{}CREA-\allowbreak{}IESR-\allowbreak{}2023} & terverifikasi \\
Biaya kesehatan per kematian & 704.762 US\$/kematian & \texttt{P-\allowbreak{}CREA-\allowbreak{}IESR-\allowbreak{}2023} & terverifikasi \\
PLTS atap terpasang nasional Juli 2025 & 538 MWp & \texttt{B-\allowbreak{}KONTAN-\allowbreak{}ROOFTOP-\allowbreak{}2025} & terverifikasi \\
Porsi PLTS atap di Jawa-Bali & 0,8 fraksi & \texttt{assumption} & asumsi \\
Tambahan PLTS atap Jawa-Bali dengan aturan sekarang & 240 MWp/tahun & \texttt{B-\allowbreak{}KONTAN-\allowbreak{}ROOFTOP-\allowbreak{}2025} & asumsi \\
Kuota PLTS atap Jawa-Bali per tahun & 800 MW/tahun & \texttt{B-\allowbreak{}PVMAG-\allowbreak{}QUOTA} & asumsi \\
Koefisien kurva pangsa pasar vs balik modal & 0,3 1/tahun & \texttt{P-\allowbreak{}NREL-\allowbreak{}SOLARDS} & perlu cek \\
Umur PLTS atap & 25 tahun & \texttt{assumption} & asumsi \\
Lapangan kerja PLTS atap & 24,72 job-year/MW & \texttt{P-\allowbreak{}CEEW-\allowbreak{}2017} & terverifikasi \\
Keluaran puncak AC per kWp & 0,75 fraksi & \texttt{assumption} & asumsi \\
PLTS atap: Rumah tangga non subsidi & Rp15 jt/kWp, pakai sendiri 40\%, tarif Rp1.444,70 & \texttt{B-\allowbreak{}KOMPASTV-\allowbreak{}TARIF-\allowbreak{}2026} & asumsi \\
PLTS atap: Bisnis & Rp13 jt/kWp, pakai sendiri 80\%, tarif Rp1.224,83 & \texttt{D-\allowbreak{}PLN-\allowbreak{}STAT-\allowbreak{}2025} & asumsi \\
PLTS atap: Industri & Rp11 jt/kWp, pakai sendiri 90\%, tarif Rp1.085,39 & \texttt{D-\allowbreak{}PLN-\allowbreak{}STAT-\allowbreak{}2025} & asumsi \\
PLTS atap: Gedung pemerintah & Rp13 jt/kWp, pakai sendiri 80\%, tarif Rp1.583,12 & \texttt{D-\allowbreak{}PLN-\allowbreak{}STAT-\allowbreak{}2025} & asumsi \\
PLTS atap: Program rumah tangga subsidi & Rp12 jt/kWp, pakai sendiri 50\%, tarif Rp605,00 & \texttt{B-\allowbreak{}KOMPASTV-\allowbreak{}TARIF-\allowbreak{}2026} & asumsi \\
Lama respons pasar PLTS atap terhadap perubahan aturan & 4 tahun & \texttt{assumption} & asumsi \\
Subsidi listrik 2026 & 101.170 miliar Rp/tahun & \texttt{B-\allowbreak{}KONTAN-\allowbreak{}SUBSIDI-\allowbreak{}2026} & terverifikasi \\
Pelanggan subsidi 2026 & 44.880.000 pelanggan & \texttt{B-\allowbreak{}KONTAN-\allowbreak{}SUBSIDI-\allowbreak{}2026} & terverifikasi \\
Porsi subsidi di Jawa-Bali & 0,61 fraksi & \texttt{D-\allowbreak{}PLN-\allowbreak{}STAT-\allowbreak{}2025} & asumsi \\
\bottomrule
\end{longtable}

}

\section{Daftar Unit PLTU}
\label{app:data}

Unit PLTU operasi terhubung jaringan PLN di Jawa-Bali \parencite{D-GEM-GCPT}. Tahun pensiun baseline memakai rencana pensiun GEM jika ada, akhir PPA 30 tahun untuk IPP, atau umur desain 40 tahun untuk PLTU milik PLN.

{\small
\begin{longtable}{p{0.22\textwidth} p{0.15\textwidth} r r l r r}
\toprule
PLTU & Unit & MW & Mulai & Tipe & Pensiun baseline & tCO$_2$/MWh \\
\midrule
\endhead
Adipala & Unit 1 & 660 & 2015 & PLN & 2055 & 0,94 \\
Batang (Central Java Power) & Unit 1 & 950 & 2022 & IPP & 2052 & 0,84 \\
Batang (Central Java Power) & Unit 2 & 950 & 2022 & IPP & 2052 & 0,84 \\
Celukan Bawang & Unit 1 & 142 & 2015 & IPP & 2045 & 1,15 \\
Celukan Bawang & Unit 2 & 142 & 2015 & IPP & 2045 & 1,15 \\
Celukan Bawang & Unit 3 & 142 & 2015 & IPP & 2045 & 1,15 \\
Cilacap Sumber & Phase I Unit 1 & 300 & 2006 & IPP & 2036 & 1,15 \\
Cilacap Sumber & Phase I Unit 2 & 300 & 2006 & IPP & 2036 & 1,15 \\
Cilacap Sumber & Phase II Unit 3 & 660 & 2016 & IPP & 2046 & 0,94 \\
Cilacap Sumber & Phase III Unit 4 & 1.000 & 2019 & IPP & 2049 & 0,84 \\
Cirebon & Unit 1 & 660 & 2012 & IPP & 2042 & 0,94 \\
Cirebon & Unit 2 & 924 & 2023 & IPP & 2053 & 0,84 \\
Indramayu & Unit 1 & 330 & 2010 & PLN & 2050 & 1,06 \\
Indramayu & Unit 2 & 330 & 2011 & PLN & 2051 & 1,06 \\
Indramayu & Unit 3 & 330 & 2011 & PLN & 2051 & 1,06 \\
Jawa-7 & Unit 1 & 991 & 2019 & IPP & 2049 & 0,84 \\
Jawa-7 & Unit 2 & 991 & 2020 & IPP & 2050 & 0,84 \\
Labuan & Unit 1 & 300 & 2009 & PLN & 2049 & 1,06 \\
Labuan & Unit 2 & 300 & 2010 & PLN & 2050 & 1,06 \\
Lontar & Unit 1 & 315 & 2011 & PLN & 2051 & 1,06 \\
Lontar & Unit 2 & 315 & 2012 & PLN & 2052 & 1,06 \\
Lontar & Unit 3 & 315 & 2012 & PLN & 2052 & 1,06 \\
Lontar & Unit 4 & 315 & 2024 & PLN & 2064 & 0,92 \\
Pacitan & Unit 1 & 315 & 2011 & PLN & 2051 & 1,06 \\
Pacitan & Unit 2 & 315 & 2011 & PLN & 2051 & 1,06 \\
Paiton & Unit 1 & 400 & 1994 & PLN & 2030 & 1,19 \\
Paiton & Unit 2 & 400 & 1993 & PLN & 2033 & 1,19 \\
Paiton Baru & Unit 09 & 660 & 2012 & PLN & 2052 & 0,97 \\
Paiton-1 & Unit 7 & 615 & 1999 & IPP & 2029 & 1,06 \\
Paiton-1 & Unit 8 & 615 & 1999 & IPP & 2029 & 1,06 \\
Paiton-2 & Unit 1 & 660 & 2000 & IPP & 2030 & 1,06 \\
Paiton-2 & Unit 2 & 660 & 2000 & IPP & 2030 & 1,06 \\
Paiton-3 & Unit 1 & 815 & 2012 & IPP & 2042 & 0,94 \\
Pelabuhan Ratu & Unit 1 & 350 & 2013 & PLN & 2045 & 1,06 \\
Pelabuhan Ratu & Unit 2 & 350 & 2013 & PLN & 2045 & 1,06 \\
Pelabuhan Ratu & Unit 3 & 350 & 2013 & PLN & 2045 & 1,06 \\
Rembang & Unit 1 & 315 & 2011 & PLN & 2051 & 1,06 \\
Rembang & Unit 2 & 315 & 2011 & PLN & 2051 & 1,06 \\
Serang & Unit 1 & 660 & 2017 & IPP & 2047 & 0,85 \\
Suralaya & Phase I Unit 1 & 400 & 1984 & PLN & 2030 & 1,19 \\
Suralaya & Phase I Unit 2 & 400 & 1985 & PLN & 2030 & 1,19 \\
Suralaya & Phase I Unit 3 & 400 & 1988 & PLN & 2030 & 1,19 \\
Suralaya & Phase I Unit 4 & 400 & 1989 & PLN & 2030 & 1,19 \\
Suralaya & Phase II Unit 1 & 600 & 1996 & PLN & 2036 & 1,10 \\
Suralaya & Phase II Unit 2 & 600 & 1997 & PLN & 2037 & 1,06 \\
Suralaya & Phase II Unit 3 & 600 & 1997 & PLN & 2037 & 1,06 \\
Suralaya & Unit 08 & 625 & 2011 & PLN & 2051 & 0,97 \\
Suralaya & Unit 09 & 1.000 & 2025 & IPP & 2055 & 0,84 \\
Suralaya & Unit 10 & 1.000 & 2025 & IPP & 2055 & 0,84 \\
Tanjung Awar-Awar & Unit 1 & 350 & 2012 & PLN & 2052 & 1,06 \\
Tanjung Awar-Awar & Unit 2 & 350 & 2013 & PLN & 2053 & 1,06 \\
Tanjung Jati B & Unit 1 & 660 & 2006 & IPP & 2036 & 1,04 \\
Tanjung Jati B & Unit 2 & 660 & 2006 & IPP & 2036 & 1,04 \\
Tanjung Jati B & Unit 3 & 660 & 2011 & IPP & 2041 & 0,95 \\
Tanjung Jati B & Unit 4 & 660 & 2012 & IPP & 2042 & 0,95 \\
Tanjung Jati B & Unit 5 & 1.000 & 2022 & IPP & 2047 & 0,84 \\
Tanjung Jati B & Unit 6 & 1.000 & 2022 & IPP & 2047 & 0,84 \\
\bottomrule
\end{longtable}

}
